% =============================================================================
%                              CORE COMMANDS
% =============================================================================
%
% Shared LaTeX3 helpers used by the package entry point and all plugins.
%
% This module defines Prime/Index/Exponent wrapping, split paired delimiters,
% lightweight command templates, and a fallback braced-operator declaration.

\ExplSyntaxOn

% =============================================================================
%                          PRIME INDEX EXPONENT HELPERS
% =============================================================================

% -----------------------------------------------------------------------------
% Public Commands
% -----------------------------------------------------------------------------
%
% \piefication
%
% Purpose:
%   Replaces an existing command with a wrapper that absorbs trailing prime,
%   index, and exponent tokens through mathcommand's PIE parser.
%
% Parameters:
%   #1 - Existing command control sequence to wrap.
%
% Usage:
%   \newmathcommand{\Foo}{\mathcal{F}}
%   \piefication{\Foo}
%   \Foo'_i^2

\NewDocumentCommand{\piefication}{m}{
    \exp_args:Nc \NewCommandCopy {old_\cs_to_str:N #1} {#1}
    \cs_set:cpn {internal_\cs_to_str:N #1:w} {
        \exp_args:Nc \__mathcommand_absorb_PIE:nw {old_\cs_to_str:N #1}
    }
    \exp_args:Nnc \RenewCommandCopy {#1} {internal_\cs_to_str:N #1:w}
}

% =============================================================================
%                            SPLIT DELIMITER HELPERS
% =============================================================================

% -----------------------------------------------------------------------------
% Public Commands
% -----------------------------------------------------------------------------
%
% \DeclarePairedSplitDelimiter
%
% Purpose:
%   Declares a paired delimiter command that can split its body at a vertical
%   bar. This is intended for set-builder notation, conditional probabilities,
%   and similar expressions.
%
% Parameters:
%   #1 - Public command to create.
%   #2 - Left delimiter token.
%   #3 - Right delimiter token.
%
% Generated command parameters:
%   *  - Uses the configured starred delimiter behavior.
%   [] - Optional delimiter size command, such as \big or \Big.
%   #1 - Body to typeset. A single | separates left and right parts.
%
% Usage:
%   \DeclarePairedSplitDelimiter{\set}{\{}{\}}
%   \set{x | x > 0}
%   \set[\Big]{x | x^2 > 1}

\NewDocumentCommand{\DeclarePairedSplitDelimiter}{mmm}{
    \exp_args:Nc\DeclarePairedDelimiter{internal_base_\cs_to_str:N#1}{#2}{#3}
    \exp_args:Nc\DeclarePairedDelimiterX{internal_split_\cs_to_str:N#1}[2]{#2}{#3}
    {##1 \nonscript\:\delimsize\vert\allowbreak\nonscript\:\mathopen{} ##2}
    \exp_args:Nc\NewDocumentCommand{internal_\cs_to_str:N#1}{somm}{
        \IfNoValueTF{##2}{
            \IfNoValueTF{##4}{
              \IfBooleanTF{##1}{
                \use:c{internal_base_\cs_to_str:N#1}{##3}
              }{
                \use:c{internal_base_\cs_to_str:N#1}*{##3}
              }
            }{
              \IfBooleanTF{##1}{
                \use:c{internal_split_\cs_to_str:N#1}*{##3}{##4}
              }{
                \use:c{internal_split_\cs_to_str:N#1}{##3}{##4}
              }
            }
        }{
            \IfNoValueTF{##3}{
                \use:c{internal_base_\cs_to_str:N#1}*[##2]{##3}
            }{
                \use:c{internal_split_\cs_to_str:N#1}*[##2]{##3}{##4}
            }
        }
    }
    \NewDocumentCommand{#1}{so >{\SplitArgument{1}{|} } m}{
        \IfNoValueTF{##2}{
          \IfBooleanTF{##1}{
            \use:c{internal_\cs_to_str:N#1}*##3
          }{
            \use:c{internal_\cs_to_str:N#1}##3
          }
        }{
            \use:c{internal_\cs_to_str:N#1}[##2]##3
        }
    }
}

% =============================================================================
%                              TEMPLATE HELPERS
% =============================================================================

% -----------------------------------------------------------------------------
% Public Commands
% -----------------------------------------------------------------------------
%
% \NewTemplateCommand
%
% Purpose:
%   Defines a reusable command template whose first argument list is supplied in
%   angle brackets and whose ordinary arguments are supplied by the caller.
%
% Parameters:
%   #1 - Template command to define.
%   #2 - Template parameter signature inside angle brackets.
%   #3 - Runtime parameter signature.
%   #4 - Replacement text.
%
% Usage:
%   \NewTemplateCommand\tmpl_example<mm>{m}{...}

\NewDocumentCommand{\NewTemplateCommand}{m r<> m m}{
    \NewDocumentCommand{#1}{#2#3}{#4}
}

% -----------------------------------------------------------------------------
% Public Commands
% -----------------------------------------------------------------------------
%
% \TemplateInstance
%
% Purpose:
%   Creates a zero-argument public command by partially applying values to a
%   template command.
%
% Parameters:
%   #1 - Command to define.
%   #2 - Template command to instantiate.
%   #3 - Comma-separated template values inside angle brackets.
%
% Usage:
%   \TemplateInstance\foo\tmpl_example<a,b>

\NewDocumentCommand{\TemplateInstance}{m m >{\SplitList {,}} r<>}{
    \NewDocumentCommand{#1}{}{#2#3}
}

% -----------------------------------------------------------------------------
% Public Commands
% -----------------------------------------------------------------------------
%
% \TemplateInstancePIE
%
% Purpose:
%   Creates a template instance and immediately wraps it with \piefication so it
%   can absorb prime, index, and exponent syntax.
%
% Parameters:
%   #1 - Command to define.
%   #2 - Template command to instantiate.
%   #3 - Comma-separated template values inside angle brackets.
%
% Usage:
%   \TemplateInstancePIE\od\tmpl_derivative_frac<\dif, \frac>

\NewDocumentCommand{\TemplateInstancePIE}{m m >{\SplitList {,}} r<>}{
    \NewDocumentCommand{#1}{}{#2#3}
    \piefication{#1}
}

% =============================================================================
%                         BRACED OPERATOR FALLBACK
% =============================================================================

% -----------------------------------------------------------------------------
% Public Commands
% -----------------------------------------------------------------------------
%
% \DeclareBracedMathOperator
%
% Purpose:
%   Provides a fallback operator declaration when the braces module is not
%   loaded. The braces module renews this command with richer delimiter support.
%
% Parameters:
%   *  - Optional star requests a limit-style operator.
%   #1 - Operator command to define.
%   #2 - Operator name to typeset.
%
% Usage:
%   \DeclareBracedMathOperator{\rank}{rank}
%   \DeclareBracedMathOperator*{\argmax}{argmax}

\ProvideDocumentCommand{\DeclareBracedMathOperator}{smm}{
  \IfBooleanTF{#1}{\DeclareMathOperator*{#2}{#3}}{\DeclareMathOperator{#2}{#3}}
}

\ExplSyntaxOff
